\documentclass[11pt]{article}
\usepackage[a4paper,margin=25mm]{geometry}
\usepackage{amsmath,amssymb,amsthm}
\newtheorem{lemma}{Lemma}
\title{Generation by panels compatible with a partial family}
\author{Human Compression project}\date{}
\begin{document}\maketitle
Let $\mathcal P_B$ be the acquired complete panel set for a fixed
$B\leq G$. Write $P\sim Q$ when the acquired rank-two test accepts the
distinct panels $P,Q$. Every family under consideration has mutually
compatible panels. This is a necessary condition; no converse is used.
For a selected family $T\subseteq\mathcal P_B$, put
\[
 C(T)=T\cup\{P\in\mathcal P_B\setminus T:
                 P\sim Q\text{ for every }Q\in T\},
 \qquad J(T)=\langle B,P\mid P\in C(T)\rangle.
\]
\begin{lemma}
If $J(T)<G$, no admissible family containing $T$ generates $G$.
The exclusion is invariant under conjugation by $N_G(B)$ and under
reordering the selected panels.
\end{lemma}
\begin{proof}
Any additional panel in an admissible family must be compatible with
each member of $T$, and hence lies in $C(T)$. The whole family therefore
generates a subgroup of $J(T)$. Properness proves the first assertion.
The acquired rank-two conditions are properties of actual subgroup
intersections, joins, maximality and coset incidence; conjugation
preserves each of them. Thus conjugation sends $C(T)$ onto the
compatible set of the conjugate family, and sends $J(T)$ to its
conjugate. Reordering $T$ changes neither set.
\end{proof}

Consequently this exclusion can be inserted into the acquired maximal
augmentation search for families generating $G$. For any such final
family $F$, the acquired augmentation theorem supplies a chain of
admissible prefixes leading to a conjugate and reordering of $F$.
For each prefix $T$ on that chain, $F\subseteq C(T)$ and therefore
$J(T)=G$. No required chain is removed. Equivariance ensures that
replacing a prefix by the representative already chosen by the
normalizer and type action does not alter this conclusion. The
original acceptance tests and the original equivalence test still
apply; this additional necessary condition accepts no new family.

One may start from an already complete saved level of the old search.
Such a level may contain prefixes subsequently ruled out by this bound.
They can be removed or skipped by the same argument. Earlier complete
full-group families remain valid. Later level sizes describe the search
with this exclusion, rather than the enumeration of all proper generated
families. The distinction is retained in the execution records.

The computation forms the symmetric compatibility matrix using the
unchanged rank-two test, then forms each actual compatible panel join.
For a rejected prefix it records that prefix, its complete compatible
index set and explicit generators for the proper join. Its correctness
depends on the acquired complete panel set, complete earlier level when
resuming, and exact finite-group operations. This is the previous
generation bound applied after a necessary local restriction; no new
classification or whole-source claim is made by the bound alone.
\end{document}
